\section*{Abstract}
We study the microstructure of Polymarket, the largest on-chain
prediction market, using a continuous tick-level archive of the
public WebSocket order-book feed (30 billion events over 52 days)
joined to the authoritative on-chain trade record. On a pre-registered
stratified panel of 600 markets we report eight stylized facts: a
longshot spread premium; a depth-concentration profile closer to a
uniform geometric grid than to the top-of-book pattern often assumed
for prediction markets; a null block-clock alignment effect; broad
maker-wallet diversity with a concentrated tail; category-conditional
differences in effective spread; a sub-50 ms median archive-ingestion
delay with a multi-second tail; a self-counterparty wash share with
median 1\% and a 22\% upper tail, where the comparison to the
network-classifier benchmarks of \citet{cong-li-tang-yang-2023} on
unregulated cryptocurrency token exchanges is a sanity bound rather
than an apples-to-apples reference because the venues face different
wash incentives; and a
depth-decay-near-resolution effect with a within-category slope of
0.55 on log seconds-to-close ($t=3.85$). The paper also contributes a
measurement result: trade direction inferred from Polymarket's public
order-book feed agrees with on-chain ground truth only $\sim 59\%$
of the time (volume-weighted across the top-100 stratum and four
disjoint 7-day windows; panel mean 0.615, market-clustered bootstrap
95\% CI $[0.58, 0.65]$, IQR $[0.53, 0.68]$),
barely above the 50\% chance baseline that direction-dependent
microstructure measures need. On the comparable subset of the
top-100 panel, the effective half-spread changes sign between feed-
and on-chain trade directions on 67\% of markets in the first 7-day
window (50\% in a second non-overlapping window), and Kyle's
$\lambda$ changes sign on 60\% (43\% in the second window); both
windows fail to recover the on-chain sign at anything close to the
$\sim 80\%$ rate that Lee-Ready achieves on equity venues.
Microstructure work on Polymarket therefore
needs to source trade direction from on-chain \texttt{OrderFilled}
events; we release a replication package that performs the join.
